\subsection{The nontrivial primitive-cover counterexample}

\begin{claim}{B.R5}{Refuted}{Universal nontrivial primitive-zero cyclic covers}
The assertion that every genus-two curve in characteristic two has
a connected cyclic \emph{\'etale} cover of odd degree $n>1$ whose
primitive Jacobian part has two-rank zero is false. A completely
specified counterexample is the smooth proper model over
$k=\mathbf F_2[t]/(M(t))$ of
\[
\begin{split}
 M(Z)&=Z^{38}+Z^{14}+Z^{10}+Z^9+Z^8+Z^5+Z^2+Z+1,\\
 d&=t^3+1,\qquad
 C:\ y^2+y=A(u)^3,\\
 A(u)&=(u^2+t^2u+t)/d,\qquad
 B(u)=(tu^2+u+t^2)/d .
\end{split}
\]
The polynomial $M$ is irreducible, and $C$ is supersingular of
genus two with a unique rational point $P$ above infinity.
For every such cover over $\bar k$, the primitive part has positive
two-rank. This independently confirms the particular counterexample
AS-A.86. It does not refute the version allowing degree one.
\end{claim}
\begin{proof}
We supply the geometry, the finite cutoff, and the exact certificates.

\emph{The curve and its Jacobian.}
Successive squaring modulo $M$ gives
$Z^{2^{38}}\equiv Z$ and
$\gcd(M,Z^{2^2}-Z)=\gcd(M,Z^{2^{19}}-Z)=1$.
These are precisely the irreducibility tests for degree 38; every
residue is retained in \path{code/primitive_certificates/manifest.json}.
In particular $t\notin\mathbf F_4$ and $d\ne0$.
After Artin--Schreier changes removing even powers, $A(u)^3$ has
highest odd term $(t^2/d^3)u^5$, with nonzero coefficient.
The degree-two map to $\mathbf P^1_u$ is ramified only at infinity,
with different exponent six. Riemann--Hurwitz gives genus two,
one rational point $P$ at infinity, and canonical divisor $2P$.
The Artin--Schreier $p$-rank formula, or its elementary one-pole
Cartier calculation, gives two-rank zero. In this case a basis of
regular differentials is $du,u\,du$ after the reduced quintic
equation; Cartier sends $du$ to zero and $u\,du$ to $du$, so its
stable rank is zero.

On $E:y^2+y=x^3$ there are two morphisms
\[
 (x_1,y_1)=(A(u),y),\qquad
 (x_2,y_2)=(B(u),y+(u^3+1)/d).
\]
The polynomial identity
\[
 (u^2+t^2u+t)^3+(tu^2+u+t^2)^3
       =d(u^3+1)^2+d^2(u^3+1)
\]
checks the second equation. Both maps send $P$ to the origin.
The joint map is birational onto its image, since
$u=x_2+t x_1$ and $y=y_1$. Its image therefore has normalization
of genus two and cannot lie in an elliptic subvariety or its
translate. The induced map $J_C\to E^2$ is consequently an isogeny
defined over $k$. On $E$, the chord-and-tangent formula gives
$[2](x,y)=(x^4,y^4+1)$; hence its $\mathbf F_4$-Frobenius is $[-2]$.
The $\mathbf F_{2^{38}}$-Frobenius of $E^2$, and therefore that of
$J_C$, is $[-2^{19}]$. For the last implication compose
$F_{J_C}+[2^{19}]$ with the isogeny: its connected image would lie
in a finite kernel, so it is zero. This proves supersingularity
and the exact scalar Frobenius, without assuming a torsion-group
isomorphism from an unspecified isogeny.

\emph{The theta-orbit criterion.}
A cyclic cover of odd degree $n$ corresponds to a line bundle $L$
of exact order $n$ with a trivialization of $L^n$. Its coherent
cohomology decomposes as
$\bigoplus_{i\in\mathbf Z/n}H^1(C,L^i)$, and absolute Frobenius
sends the $i$-summand semilinearly to the $2i$-summand. The
nontrivial summands have dimension one by Riemann--Roch.
For a multiplication-by-two orbit $O$ of primitive indices, its
contribution to the two-rank is $|O|$ times the eventual rank of
the orbit composite. This follows from the Artin--Schreier
sequence: the Frobenius-stable subspace has dimension equal to
$\dim_{\mathbf F_2}H^1_{\et}(D,\mathbf F_2)$. The idempotent
primitive character projector, or the image of
$\prod_{\ell\mid n}(1-\sigma^{n/\ell})$, selects exactly these
primitive summands.

Let $\mathcal B$ be the cokernel of
$\mathcal O_{C^{(1)}}\to F_*\mathcal O_C$. In characteristic two it
is a line bundle with square the canonical bundle: locally write
a change of uniformizer as $v=a^2+u b^2$; its class in the
cokernel changes by $b$, whose square is $dv/du$. The cohomology
sequence tensored with $L^i$ identifies the kernel of the
Frobenius arrow with $H^0(\mathcal B\otimes L^i)$.
There is just one geometric theta characteristic because
$J_C[2](\bar k)=0$, so $\mathcal B=\mathcal O(P)$.
The corresponding theta locus is precisely the symmetric Abel
curve $\Theta=C-P$. Thus zero primitive two-rank requires at
least one theta point in every multiplication-by-two orbit of
primitive points of the cyclic subgroup $\langle L\rangle$.

\emph{A complete finite bound.}
Write $f=\operatorname{ord}_n(2)$,
$e=\gcd(f,19)$, and $M_n=\varphi(n)/e$.
Scalar Frobenius and symmetry of $\Theta$ show that one point
in a required orbit brings its entire multiplication-by-$2^{19}$
orbit into $\Theta$. Consequently there are at least $M_n$
primitive theta points, and also zero. For a nonzero Jacobian
difference, at most two ordered pairs of distinct points of
$\Theta$ realize it: two representations yield two linearly
equivalent degree-two divisors, and the unique hyperelliptic
pencil gives the sole alternative representation. Hence
\[
 M_n(M_n+1)\le2(n-1).
\]
For odd $n$, prime-power multiplication proves
$\varphi(n)\ge(2/3)n^{5/6}$. Indeed
$(\ell-1)^6\ge\ell^5$ for every prime $\ell\ge5$, and the possible
factor at three is at least $64/243>(2/3)^6$.
It follows that $n<10\cdot19^3=68590$.

The standard-library program \path{code/primitive_conductors.py}
exhaustively enumerates all odd integers in that interval and gives
\[
191,229,361,419,457,571,573,647,687,955,1083,1257,1371.
\]
It then imposes the full difference constraint on the mandatory
packets $H=\langle 2^{19},-1\rangle\subset(\mathbf Z/n)^\times$.
Already $\{0\}\cup H$ excludes all but $191,229,457,1257$.
For $457$, every coset of $U=\langle2,-1\rangle$ must contain an
$H$-packet. There are six such cosets. Normalize the first packet
to $H$ by changing a cyclic generator, and enumerate every
remaining packet. The surviving search depths are
$1,17,226,879,262,0$, so there is no admissible selection.
For $1257$, there are two $U$-cosets and each mandatory packet
has 22 elements, half in each nonzero residue class modulo three.
With zero, the residue counts are $(1,22,22)$, forcing
$924$ ordered distinct differences divisible by three. Only
$2(1257/3-1)=836$ are permitted. Thus only $191,229$ remain.
The script retains every input, packet, candidate list and count.

\emph{The two full polynomial certificates.}
For $E$ put $r_n(X)=x([n]R)$ when $X=x(R)$. The addition formula gives
\[
 r_1=X,\qquad r_{2m}=r_m^4,\qquad
 r_{2m+1}=X+(r_m+r_{m+1})^{-2}.
\]
Compute these as reduced rational functions over $\mathbf F_2$.
Independently form
\[
 \psi_1=\psi_2=\psi_4=1,\quad\psi_3=X^4+X,\quad
 \psi_{2m}=\psi_m^4,\quad
 \psi_{2m+1}=\psi_{m+2}\psi_m^3+\psi_{m-1}\psi_{m+1}^3.
\]
For each required $n$ the program verifies that the reduced
denominator of $r_n$ is exactly $\psi_n^2$, with coprime numerator.
Thus its zeros are exactly the finite coordinates of nonzero
$n$-torsion, without importing a division-polynomial formula
that excludes characteristic two.
\[
\begin{array}{c|r|r|r}
 n&\deg\psi_n&\deg\operatorname{num}r_n&
                  \deg\operatorname{den}r_n\\ \hline
191&18240&36481&36480\\
229&26220&52441&52440
\end{array}
\]
Exact arithmetic in the stated field gives
\[
 U_n(u)\psi_n(A(u))+V_n(u)\psi_n(B(u))=1
                    \quad(n=191,229).
\]
The composed polynomial degrees are respectively $36480,52440$;
both B\'ezout coefficients have degree four less in each case.
All eight polynomials are supplied as complete binary coefficient
arrays, with ascending unsigned little-endian 64-bit coefficients;
bit $i$ represents $t^i$. The manifest gives SHA-256 hashes.
The producer checks both whole identities. The separately written
\path{code/primitive_verify.py} loads the arrays, verifies their
hashes and degrees, independently composes the polynomials over
$\mathbf F_2$, and verifies the two whole identities again.

If a finite point $Q$ of $C$ gave a nonzero Abel class killed by
one of these integers, both of its finite elliptic images would
be killed by the integer. Its coordinate $u$ would therefore be a
common zero of the two composed polynomials, contradicting the
B\'ezout identity. The point at infinity is the origin and cannot
be primitive. Every possible conductor has now been excluded,
over $\bar k$ itself. This proves the counterexample.
\end{proof}

\begin{claim}{B.R6}{Open}{The literal degree-one-permitted primitive-cover assertion}
For every proper curve of genus at least two over
$\overline{\mathbf F}_p$, does there exist a connected cyclic
\emph{\'etale} cover of odd degree prime to $p$, allowing degree one,
whose primitive part has $p$-rank zero, with the convention
$P_1=J_C$?
\end{claim}
The preceding counterexample does not settle this statement.
The identity cover proves it for every curve of $p$-rank zero,
including that example. For positive $p$-rank, degree one fails.
The exact Frobenius-orbit formulation above extends to arbitrary
genus with character spaces of dimension $g-1$: one needs
nilpotence of each orbit composite, not just a noninvertible
arrow. The conductor bound proved here uses scalar Frobenius
and the one-dimensional genus-two character spaces, so it does
not bound the positive-$p$-rank cases. We tried transferring the
scalar obstruction, using a single nonordinary character, and
using the identity convention; each fails for this stated reason.
The smallest unresolved scope in this work is positive-two-rank
genus-two curves in characteristic two. No such curve is claimed
to be a counterexample.

